\chapterbackmatter{Solutions to Weinberg Exercises}
\ExerciseEditorialNote

\begin{enumerate}
  \WeinbergSolution{1}
  Put \(x_a=t_a\), \(a=1,2\), and write
  \[
    \gamma(\boldsymbol\xi)
      =\exp(i\xi_at_a),\qquad
    r=(\xi_1^2+\xi_2^2)^{1/2},\qquad
    n_a=\xi_a/r.
  \]
  We use \([t_i,t_j]=i\epsilon_{ijk}t_k\) and
  \(\epsilon_{12}=+1\).  For an infinitesimal transformation
  \(g=1+i\epsilon_at_a+i\epsilon_3t_3\), the factorization
  \(g\gamma(\xi)=\gamma(\xi+\delta\xi)\exp(i\vartheta t_3)\)
  gives
  \[
    \boxed{\delta\xi_a
      =r\cot r\,\epsilon_a
       +(1-r\cot r)n_a n_b\epsilon_b
       +\epsilon_3\epsilon_{ab}\xi_b,}
  \]
  \[
    \vartheta=\epsilon_3+
      \frac{\tan(r/2)}{r}
      (\xi_1\epsilon_2-\xi_2\epsilon_1).
  \]
  The limit at \(r=0\) is smooth.  In particular,
  \(\delta\xi_a=\epsilon_a+\epsilon_3\epsilon_{ab}\xi_b
  +O(\xi^2\epsilon)\).

  The Maurer--Cartan form is
  \[
    \gamma^{-1}\partial_\mu\gamma
      =it_aD_{a\mu}+it_3E_\mu,
  \]
  where
  \[
    \boxed{D_{a\mu}
      =\frac{\sin r}{r}\partial_\mu\xi_a
       +\left(1-\frac{\sin r}{r}\right)
       n_a n_b\partial_\mu\xi_b,}
  \]
  \[
    \boxed{E_\mu
      =\frac{1-\cos r}{r^2}
       (\xi_1\partial_\mu\xi_2-\xi_2\partial_\mu\xi_1).}
  \]
  Thus a field \(\chi_q\) on which \(t_3\) has eigenvalue \(q\)
  transforms as
  \(\chi_q\mapsto e^{iq\vartheta}\chi_q\), and its covariant derivative is
  \[
    \boxed{\mathcal D_\mu\chi_q
      =(\partial_\mu+iqE_\mu)\chi_q.}
  \]

  The broken generators form an irreducible real doublet of \(SO(2)\).
  Hence the unique Goldstone Lagrangian with two derivatives is
  \[
    \mathcal L_2=-\frac{F^2}{2}D_{a\mu}D_a^\mu.
  \]
  With canonically normalized fields \(\pi_a=F\xi_a\), it expands as
  \[
    \mathcal L_2=-\frac12\partial_\mu\pi_a\partial^\mu\pi_a
      +\frac{1}{6F^2}
       \left[
        \boldsymbol\pi^2(\partial_\mu\boldsymbol\pi)^2
        -(\boldsymbol\pi\cdot\partial_\mu\boldsymbol\pi)^2
       \right]+\cdots.
  \]
  The on-shell invariant amplitude is therefore
  \[
    \boxed{\mathcal M_{ab,cd}
      =\frac1{F^2}\left(
       \delta_{ab}\delta_{cd}s+
       \delta_{ac}\delta_{bd}t+
       \delta_{ad}\delta_{bc}u\right)+O(Q^4).}
  \]

  The rule for matter is most compactly stated as a charge-selection
  rule.  Define \(D_\mu^\pm=(D_{1\mu}\mp iD_{2\mu})/\sqrt2\); these have
  \(SO(2)\) charges \(\pm1\).  The most general bilinear Lagrangian with
  at most one derivative is the most general Lorentz scalar made from
  \[
    \overline\chi_{q'}\chi_q,\qquad
    \overline\chi_{q'}\mathcal D_\mu\chi_q,\qquad
    D_\mu^\pm\overline\chi_{q'}\Gamma^\mu\chi_q,
  \]
  with respectively \(q'=q\), \(q'=q\), and
  \(q'=q\pm1\), and with independent coefficients for every allowed
  Lorentz structure \(\Gamma^\mu\).  For a single species of fixed
  charge, this reduces to its arbitrary invariant mass and covariant
  kinetic terms; a term linear in \(D^\pm_\mu\) requires another species
  whose charge differs by one.

  Finally, the three-vector obtained by rotating the unbroken axis has
  third component
  \[
    V_3(\xi)=V\cos r.
  \]
  Up to an overall constant this is the only derivative-free Goldstone
  function with the required transformation law.  Choosing
  \[
    \Delta\mathcal L=F^2m^2(\cos r-1)
      =-\frac12m^2\boldsymbol\pi^2
       +\frac{m^2}{24F^2}(\boldsymbol\pi^2)^2+\cdots
  \]
  gives both Goldstones mass \(m\).  Adding the contact term from this
  potential changes the amplitude to
  \[
    \boxed{\mathcal M_{ab,cd}
      =\frac1{F^2}\left[
       \delta_{ab}\delta_{cd}(s-m^2)
       +\delta_{ac}\delta_{bd}(t-m^2)
       +\delta_{ad}\delta_{bc}(u-m^2)\right]+O(Q^4).}
  \]

  \WeinbergSolution{2}
  A vacuum for the breaking \(SU(N)\to SU(N-1)\) may be represented by a
  unit vector \(v\in\mathbb C^N\); its stabilizer is the \(SU(N-1)\) that
  fixes \(v\).  If the defining-representation perturbation is represented
  by a fixed vector \(c\), the first-order vacuum energy has the form
  \[
    E_1(v)=-\kappa(c^\dagger v+v^\dagger c).
  \]
  A phase and the sign of \(v\) can be chosen so that the minimum has
  \(v=c/\lVert c\rVert\).  The source and the aligned vacuum then have the
  same stabilizer:
  \[
    \boxed{H_{\rm exact}\cap H_{\rm vac}=SU(N-1).}
  \]

  An adjoint perturbation is a traceless Hermitian matrix \(A\).  The
  \(SU(N-1)\)-invariance of the unperturbed vacuum implies that its
  first-order energy can depend on \(v\) only through
  \[
    E_1(v)=\kappa\,v^\dagger A v+\text{constant}.
  \]
  The Rayleigh--Ritz theorem aligns \(v\) with an eigenvector belonging to
  the lowest or highest eigenvalue, according to the sign of \(\kappa\).
  The completely unbroken group is consequently
  \[
    \boxed{H_{\rm complete}=C_{SU(N)}(A)\cap SU(N-1)_v,}
  \]
  where \(C_{SU(N)}(A)\) is the centralizer of \(A\).

  This formula also displays a qualification that is essential for the
  adjoint representation: the answer depends on eigenvalue degeneracies.
  If the selected eigenvalue has multiplicity \(k\), while the other
  distinct eigenvalues have multiplicities \(n_2,n_3,\ldots\), then
  \[
    H_{\rm complete}
      =S\!\left(U(k-1)\times U(n_2)\times U(n_3)\times\cdots\right).
  \]
  For a generic adjoint source all eigenvalues are distinct, and hence
  \[
    \boxed{H_{\rm complete}=U(1)^{N-2}.}
  \]
  Enhanced answers follow at degenerate points.  For example, a source
  with multiplicities \(1\) and \(N-1\) can leave \(SU(N-1)\), or
  \(S(U(N-2)\times U(1))\), depending on which eigenspace is selected by
  the alignment.

  \WeinbergSolution{3}
  It is convenient to put \(f=F_\pi/2\), so that \(f\simeq92\,\mathrm{MeV}\)
  is the more common chiral-perturbation-theory normalization.  Isospin
  and crossing give
  \[
    \mathcal M_{ab,cd}
      =\delta_{ab}\delta_{cd}A(s,t,u)
       +\delta_{ac}\delta_{bd}A(t,s,u)
       +\delta_{ad}\delta_{bc}A(u,t,s).
  \]
  Through one loop, expressed in terms of the physical \(m_\pi\) and
  \(f\), the finite answer may be written
  \[
    A(s,t,u)=\frac{s-m_\pi^2}{f^2}+B(s,t,u)+C(s,t,u)+O(Q^6),
  \]
  with \(s+t+u=4m_\pi^2\) and
  \[
  \begin{split}
    B(s,t,u)=\frac1{6f^4}\Big\{&
      3(s^2-m_\pi^4)\overline J(s)\\
      &+[t(t-u)-2m_\pi^2t+4m_\pi^2u-2m_\pi^4]\overline J(t)\\
      &+[u(u-t)-2m_\pi^2u+4m_\pi^2t-2m_\pi^4]\overline J(u)\Big\},
  \end{split}
  \]
  \[
  \begin{split}
    C(s,t,u)=\frac1{96\pi^2f^4}\Big\{&
      2\left(\overline\ell_1-\frac43\right)(s-2m_\pi^2)^2\\
      &+\left(\overline\ell_2-\frac56\right)
       [s^2+(t-u)^2]
       +12m_\pi^2s(\overline\ell_4-1)\\
      &-3m_\pi^4(\overline\ell_3+4\overline\ell_4-5)\Big\}.
  \end{split}
  \]
  Here
  \[
    \overline J(z)=\frac1{16\pi^2}
      \left[
       2+\sigma(z)\ln\frac{\sigma(z)-1}{\sigma(z)+1}
      \right],
    \qquad
    \sigma(z)=\sqrt{1-\frac{4m_\pi^2}{z}},
  \]
  with the logarithm continued with the Feynman prescription.  The
  scale-independent \(\overline\ell_i\) are the four renormalized
  \(SU(2)\) low-energy constants at this order.  Equivalently one may use
  scale-dependent \(\ell_i^r(\mu)\); their \(\mu\)-dependence then cancels
  that of the unsubtracted loop function.

  The first term is precisely
  \(4(s-m_\pi^2)/F_\pi^2\), as in Eq.~\eqref{eq:19.5.35}.  The function
  \(B\) contains every nonanalytic one-loop contribution and has the
  correct two-pion cuts in all three channels; \(C\) is the most general
  crossing-symmetric local polynomial at order \(Q^4\), including the
  mass and decay-constant renormalizations.  Thus the displayed result is
  finite, crossing symmetric, and independent of the subtraction scale.

  \WeinbergSolution{4}
  Let
  \(Q_5^\pm=\int d^3x\,A_0^\pm(x)\).  The current algebra is
  \[
    [Q_5^+,Q_5^-]=2I_3.
  \]
  Take its matrix element between one-\(\pi^-\) states and insert a
  complete set between the two charges.  PCAC isolates the one-soft-pion
  pole,
  \[
    \bra{\Omega}A_\mu^a(0)\ket{\pi^b(q)}
       =\frac{iF_\pi}{2}\,q_\mu\delta^{ab},
  \]
  where the factor \(1/2\) converts the chapter's
  \(F_\pi\simeq184\,\mathrm{MeV}\) to the conventional decay constant.
  The two orders of the charges become the absorptive parts of forward
  \(\pi^+\pi^-\) and \(\pi^-\pi^-\) amplitudes.  The equal-time
  commutator fixes their difference rather than either amplitude
  separately.

  A once-subtracted forward dispersion relation then gives the
  off-shell form of the Adler sum rule
  \[
    \boxed{
    1=\pi F_\pi^2
      \int_{4m_\pi^2}^{\infty}
      \frac{ds}{(s-m_\pi^2)^2}
      \left[
       \operatorname{Im}\mathcal F_{\pi^+\pi^-}(0;s,0)
       -\operatorname{Im}\mathcal F_{\pi^-\pi^-}(0;s,0)
      \right].}
  \]
  The first argument says that the incident pion has \(q^2=0\); the
  target pion remains on shell.  The normalization of
  \(\mathcal F\) is fixed by the displayed equation, or, equivalently,
  by writing the standard convention \(f_\pi=F_\pi/2\), in which the
  prefactor is \(4\pi f_\pi^2\).

  In the exact chiral limit the off-shell extrapolation disappears and
  the lower endpoint moves to zero:
  \[
    \boxed{
    1=\pi F_\pi^2
      \int_0^\infty\frac{ds}{s^2}
      \left[
       \operatorname{Im}\mathcal F_{\pi^+\pi^-}(s,0)
       -\operatorname{Im}\mathcal F_{\pi^-\pi^-}(s,0)
      \right].}
  \]
  The optical theorem converts the two absorptive parts into the
  corresponding total cross sections.  This is the pion-target analog
  of Eq.~\eqref{eq:19.5.71}: the nucleon matrix element of the same
  commutator produces \(g_A^2\), whereas the pion matrix element is fixed
  entirely by isospin and produces the number \(1\).

  \WeinbergSolution{5}
  Use generators satisfying
  \[
    [t_a,t_b]=i\epsilon_{abc}t_c,\qquad
    [t_a,x_b]=i\epsilon_{abc}x_c,\qquad
    [x_a,x_b]=i\epsilon_{abc}t_c
  \]
  and choose
  \(\gamma(\boldsymbol\xi)=\exp(i\boldsymbol\xi\cdot\boldsymbol x)\).
  Write \(r=|\boldsymbol\xi|\) and
  \(\boldsymbol n=\boldsymbol\xi/r\).  For
  \[
    g=1+i\boldsymbol\alpha\cdot\boldsymbol t
       +i\boldsymbol\beta\cdot\boldsymbol x,
  \]
  factor
  \(g\gamma(\xi)=\gamma(\xi+\delta\xi)h\), with \(h\in SU(2)_V\).
  Resolving the broken part parallel and perpendicular to
  \(\boldsymbol n\) gives
  \[
    \delta\boldsymbol\xi_{\parallel}
      =\boldsymbol\beta_{\parallel},\qquad
    \delta\boldsymbol\xi_{\perp}
      =r\cot r\,\boldsymbol\beta_{\perp},
  \]
  while an unbroken transformation rotates the coordinate vector
  linearly.  In components,
  \[
    \boxed{
    \delta\xi_a
      =(\boldsymbol\xi\times\boldsymbol\alpha)_a
       +r\cot r\,\beta_a
       +(1-r\cot r)
        \frac{\xi_a\xi_b}{r^2}\beta_b.}
  \]
  The compensating vector rotation is
  \[
    h=1+i\boldsymbol\vartheta\cdot\boldsymbol t,\qquad
    \boldsymbol\vartheta
      =\boldsymbol\alpha+
       \frac{\tan(r/2)}{r}
       \boldsymbol\xi\times\boldsymbol\beta.
  \]
  Expanding near the origin provides a useful check:
  \[
    \delta\boldsymbol\xi
      =\boldsymbol\beta+\boldsymbol\xi\times\boldsymbol\alpha
       +\frac13\boldsymbol\xi\times
        (\boldsymbol\xi\times\boldsymbol\beta)+O(\xi^4\beta).
  \]
  Thus the broken symmetry begins with a constant shift, while the
  unbroken diagonal \(SU(2)\) acts as the ordinary isovector rotation.

  \WeinbergSolution{6}
  Represent the baryon octet by the traceless matrix
  \[
    B=\begin{pmatrix}
      \Sigma^0/\sqrt2+\Lambda/\sqrt6&\Sigma^+&p\\
      \Sigma^-&-\Sigma^0/\sqrt2+\Lambda/\sqrt6&n\\
      \Xi^-&\Xi^0&-2\Lambda/\sqrt6
    \end{pmatrix},
    \qquad B\longmapsto UBU^\dagger.
  \]
  To first order in the quark-mass spurion
  \(M=\operatorname{diag}(m_u,m_d,m_s)\), the most general octet mass
  operator has only the two octet reduced matrix elements \(D,F\), plus a
  singlet:
  \[
    \delta\mathcal L
      =-b_D\operatorname{Tr}(\overline B\{M,B\})
       -b_F\operatorname{Tr}(\overline B[M,B])
       -b_0\operatorname{Tr}(M)\operatorname{Tr}(\overline BB).
  \]
  For an off-diagonal component \(B^i{}_j\), the non-singlet mass shift is
  \[
    \delta m(B^i{}_j)
      =b_D(m_i+m_j)+b_F(m_i-m_j).
  \]
  Since \(p=B^1{}_3\) and \(n=B^2{}_3\),
  \[
    \delta m_p-\delta m_n
      =(b_D+b_F)(m_u-m_d).
  \]
  Now take the isospin limit in the coefficient multiplying the much
  larger strange mass.  With
  \(\Sigma^+=B^1{}_2\) and \(\Xi^-=B^3{}_1\),
  \[
    m_\Xi-m_\Sigma=(b_D+b_F)m_s
      +O(m_u,m_d).
  \]
  The singlet coefficient \(b_0\) cancels from both differences.
  Eliminating the common reduced matrix element therefore yields
  \[
    \boxed{
    \frac{\delta m_p-\delta m_n}{m_u-m_d}
      =\frac{m_\Xi-m_\Sigma}{m_s},}
  \]
  which is Eq.~\eqref{eq:19.7.32} to the first order in quark masses
  implicit in Eq.~\eqref{eq:19.7.31}.
\end{enumerate}
